\documentclass[11pt]{article}

\usepackage[T1]{fontenc}
\usepackage{amsmath,amssymb,amsthm}
\usepackage[margin=1in]{geometry}
\usepackage[hidelinks]{hyperref}
\setlength{\emergencystretch}{3em}

\newtheorem{theorem}{Theorem}
\newtheorem{lemma}[theorem]{Lemma}
\newtheorem{proposition}[theorem]{Proposition}
\newtheorem{corollary}[theorem]{Corollary}
\newtheorem{conjecture}[theorem]{Conjecture}
\theoremstyle{remark}
\newtheorem{remark}[theorem]{Remark}

\newcommand{\R}{\mathbb R}
\newcommand{\HH}{\mathcal H}
\newcommand{\EE}{\mathcal E}
\newcommand{\one}{\mathbf 1}
\newcommand{\weakto}{\rightharpoonup}
\DeclareMathOperator{\supp}{supp}
\DeclareMathOperator{\dimH}{dim_{\mathrm H}}

\title{Exact q4 Terminal Dimension Pincer and Spectral Alternative}
\author{John Robert Lawson and Codex}
\date{Research round ending 25 July 2026}

\begin{document}
\maketitle

\begin{abstract}
This round audits the energy-equality route forced by the preceding
critical-clock theorem.  The exact first-record schedule gives the
endpoint distribution
\[
 \|u(t)\|_{L^{3,\infty}}
 \in L_t^{11/2,\infty}
 \cap\bigcap_{r>11/2}L_t^{11/2,r}.
\]
Even after granting the additional favourable surrogate assumption
\(L_t^{11/2,\infty}B^0_{3,\infty}\), the complete interpolation hull
with the Leray energy class misses every relevant
Cheskidov--Luo Besov threshold.  The refined Leslie--Shvydkoy
first-time threshold has sharp hull supremum \(3/5\).  Combining this
with the previously proved nonzero terminal energy defect shows,
conditionally and subject to external review, that its terminal
singular slice must satisfy
\[
 \frac35\leq\dimH\sigma\leq1,
 \qquad \HH^1(\sigma)=0.
\]
The pressure associated with the same clock has direct-interpolation
supercriticality at least \(8/11\), whereas the Barker--Wang theorem
could enter the dimension-closing region only below \(1/4\).
Finally, an exhaustive spectral reduction gives at least one of two
alternatives.  The defect can force infinitely many disjoint fixed-sign
nonlinear correlation events in a bounded clock band; if that alternative
fails, it escapes beyond every bounded multiple of the clock frequency and
amplifies the pointwise enstrophy.
Neither survivor is contradicted by the current budgets.  No Clay
alternative is proved.
\end{abstract}

\section{Epistemic scope and source anchors}

The solution and exact \(q=4\) schedule below are conditional.  The
endpoint-clock conversion, convex-hull optimisations, terminal
dimension pincer, pressure comparison, and spectral alternative are
repository deductions and await external mathematical review.

The imported energy-equality criteria are from Cheskidov and
Luo~\cite{CL20}.  Their arXiv v2 source states the weak-in-time theorem
at lines 194--200 and the four classical interpolation thresholds at
lines 349--357.  The imported concentration-dimension theorem is from
Leslie and Shvydkoy~\cite{LS18}; its accepted arXiv v4 source gives the
Navier--Stokes regions and function \(f\) at lines 1183--1210 and the
theorem at lines 1241--1252.  The pressure singular-set theorem is from
Barker and Wang~\cite{BW23}; its arXiv v2 source gives Theorems 1--2 at
lines 230--306.  These three results are peer reviewed.

The statement \(\HH^1(\sigma)=0\) is the classical
Caffarelli--Kohn--Nirenberg conclusion~\cite{CKN82}.  General
Leray--Hopf energy equality remains open.  Same-system adversarial
recomputation is not independent external review.  The Clay problem
remains unsolved.

\section{Conditional setting}

Let \(u\) be a smooth divergence-free finite-energy solution of
\[
 \partial_tu+u\cdot\nabla u=-\nabla\pi+\nu\Delta u,
 \qquad \nabla\cdot u=0
 \quad\hbox{on }\R^3\times[0,T^*),
\]
where \(T^*<\infty\) is its first possible singular time.  Assume a
weakly continuous Leray--Hopf continuation through \(T^*\), and write
\[
 u_*:=u(T^*).
\]

Put
\[
 M(t):=\|u(t)\|_{L^{3,\infty}}.
\]
Let \(t_j\uparrow T^*\) be exact \(q=4\) first-record times, so that
\begin{equation}
 \label{eq:q4}
 m_j:=M(t_j)\asymp2^{2j},
 \qquad
 t_{j+1}-t_j\asymp\frac{2^{-11j}}j,
 \qquad
 M(t)\leq m_j\quad(0\leq t\leq t_j).
\end{equation}
Assume the energy-efficient diffuse selected-layer hypotheses of the
critical-clock theorem.  That theorem supplies a fixed block length,
an endpoint sequence \(s_j\uparrow T^*\), and frequencies
\begin{equation}
 \label{eq:Lambda}
 \Lambda_j\asymp m_j^{7/3}j^{2/3}
\end{equation}
such that
\begin{equation}
 \label{eq:old-floor}
 \|P_{>\Lambda_j}u(s_j)\|_2\geq c_0>0.
\end{equation}
It also gives the terminal energy defect
\begin{equation}
 \label{eq:energy-defect}
 E_--\|u_*\|_2^2\geq d_0>0,
 \qquad
 E_-:=\lim_{t\uparrow T^*}\|u(t)\|_2^2.
\end{equation}
Changing \(j\) by the fixed block length changes none of the powers
used below.

\section{The exact temporal endpoint}

\begin{lemma}[Logarithmically improved endpoint distribution]
\label{lem:endpoint}
For all sufficiently large \(\lambda\),
\begin{equation}
 \label{eq:distribution}
 \bigl|\{t<T^*:M(t)>\lambda\}\bigr|
 \leq
 \frac{C\lambda^{-11/2}}{\log(e+\lambda)}.
\end{equation}
Consequently,
\begin{equation}
 \label{eq:lorentz-endpoint}
 M\in L^{11/2,\infty}(0,T^*)
 \cap
 \bigcap_{r>11/2}L^{11/2,r}(0,T^*).
\end{equation}
\end{lemma}

\begin{proof}
Choose \(j\) so that \(m_j\leq\lambda<m_{j+1}\).  The first-record
property in \eqref{eq:q4} gives
\[
 \{M>\lambda\}\subset(t_j,T^*).
\]
The geometric record tail obeys
\[
 T^*-t_j
 \lesssim
 \sum_{k\geq j}\frac{2^{-11k}}k
 \lesssim\frac{2^{-11j}}j.
\]
Since \(\lambda\asymp2^{2j}\) and
\(j\asymp\log(e+\lambda)\), this is
\eqref{eq:distribution}.

Let \(p_*=11/2\) and denote the distribution function by
\(\mu_M\).  Equation \eqref{eq:distribution} gives
\[
 \sup_{\lambda>0}\lambda^{p_*}\mu_M(\lambda)<\infty,
\]
which is the weak endpoint.  For finite \(r\), the
distribution-function formula gives
\[
 \|M\|_{L^{p_*,r}}^r
 \asymp
 \int_0^\infty
 \left[\lambda\mu_M(\lambda)^{1/p_*}\right]^r
 \frac{d\lambda}{\lambda}.
\]
The terminal tail is at most
\[
 C\int^\infty
 \frac{d\lambda}
 {\lambda[\log(e+\lambda)]^{r/p_*}},
\]
which converges when \(r>p_*\).  Smoothness controls the remaining
part.  At \(r=p_*\), this estimate is harmonic and makes no strong
\(L^{11/2}\) claim.
\end{proof}

\section{Direct Besov interpolation obstruction}

No embedding comparison between the actual endpoint
\(L_x^{3,\infty}\) and \(B^0_{3,\infty}\) is used here.  We instead grant
the additional favourable surrogate assumption
\[
 C=L_t^{11/2,\infty}B^0_{3,\infty}
\]
and combine it with
\[
 A=L_t^\infty B^0_{2,2},
 \qquad
 B=L_t^2B^1_{2,2}.
\]
Let \(a,b,c\geq0\), \(a+b+c=1\), be interpolation weights on
\(A,B,C\).  The optimistic available triple is
\begin{equation}
 \label{eq:hull-triple}
 \frac1\beta=\frac b2+\frac{2c}{11},
 \qquad
 \frac1p=\frac12-\frac c6,
 \qquad
 \alpha_{\rm av}=b.
\end{equation}

Cheskidov and Luo record the following sufficient spatial smoothness
across the four classical interpolation regions:
\begin{equation}
 \label{eq:CL-four}
 \alpha_{\rm req}=
 \begin{cases}
 \dfrac2\beta+\dfrac2p-1,
  &\beta\geq3,\ p\geq\beta,\\[1mm]
 \dfrac1\beta+\dfrac3p-1,
  &\beta\geq3,\ p\leq\beta,\\[1mm]
 \dfrac5{2\beta}+\dfrac3p-\dfrac32,
  &\beta\leq3,\ \dfrac1\beta+\dfrac2p\geq1,\\[1mm]
 \dfrac2\beta+\dfrac2p-1,
  &\beta\leq3,\ \dfrac1\beta+\dfrac2p\leq1.
 \end{cases}
\end{equation}
Their weak-in-time theorem has the same smoothness as the first line
under its strict parameter hypotheses.

\begin{proposition}[Strict four-region gap]
\label{prop:CL-gap}
No applicable point in \eqref{eq:hull-triple} reaches
\eqref{eq:CL-four}.  In the four regions, respectively,
\begin{equation}
 \label{eq:four-gaps}
 \alpha_{\rm req}-\alpha_{\rm av}
 =
 \begin{cases}
 c/33,\\[1mm]
 a/2+2c/11,\\[1mm]
 b/4-c/22,\\[1mm]
 c/33.
 \end{cases}
\end{equation}
Every applicable expression is strictly positive.
\end{proposition}

\begin{proof}
Direct substitution of \eqref{eq:hull-triple} into
\eqref{eq:CL-four} gives \eqref{eq:four-gaps}.
In the third region,
\[
 \frac1\beta+\frac2p\geq1
 \quad\Longrightarrow\quad
 \frac b2+\frac{2c}{11}\geq\frac c3
 \quad\Longrightarrow\quad
 b\geq\frac{10c}{33}.
\]
Therefore
\[
 \frac b4-\frac c{22}\geq\frac c{33}.
\]
If \(c=0\), the third-region condition forces \(b>0\), and the
original expression \(b/4\) is positive.

In the first and fourth lines, equality would require \(c=0\), hence
\(p=2\).  This contradicts \(p\geq\beta\geq3\) in the first line.  In
the fourth line, its second region inequality then forces \(b=0\),
which contradicts \(\beta\leq3\).  Equality in the second line would
require the pure \(B\) point, but that point has \(\beta=p=2\), outside
the second region.
\end{proof}

\begin{remark}
Proposition~\ref{prop:CL-gap} excludes only direct interpolation from
the stated norms.  It does not exclude a new structural estimate,
pressure cancellation, or another route to energy equality.

The calculation also covers the standard post-interpolation
Sobolev/Besov embeddings.  Spending \(h\geq0\) units of reciprocal
spatial integrability changes
\[
 \frac1p\longmapsto\frac1p-h,
 \qquad
 \alpha_{\rm av}\longmapsto\alpha_{\rm av}-3h.
\]
In the four lines of \eqref{eq:CL-four}, the gap in
\eqref{eq:four-gaps} respectively increases by \(h\), stays fixed,
stays fixed, or increases by \(h\).  Lowering the time exponent on the
finite interval only raises the required smoothness.  Thus the no-go
includes, for example, the embedding
\(L_t^2H_x^1\subset L_t^2L_x^6\).
\end{remark}

\section{Sharp Leslie--Shvydkoy hull ceiling}

Use reciprocal mixed-norm coordinates
\[
 x:=\frac1p,\qquad y:=\frac1q.
\]
The closure of the available strong interior points lies in the convex
hull of
\begin{equation}
 \label{eq:mixed-hull}
 A_0=\left(\frac12,0\right),
 \qquad
 B_0=\left(\frac16,\frac12\right),
 \qquad
 C_0=\left(\frac13,\frac2{11}\right).
\end{equation}
The endpoint \(C_0\) itself is weak in time and space, but it controls
the closure.

In Leslie--Shvydkoy Regions I--II,
\begin{equation}
 \label{eq:f12}
 f(p,q)
 =
 3-\frac{2y}{1-2x-y}.
\end{equation}
In Region III,
\begin{equation}
 \label{eq:f3}
 f(p,q)
 =
 3-
 \frac{2y(6x-1)}
 {(2-3x-3y)(1-2x)}.
\end{equation}

\begin{proposition}[Exact \(3/5\) optimisation]
\label{prop:three-fifths}
Every Leslie--Shvydkoy-admissible point in the hull
\eqref{eq:mixed-hull} satisfies
\[
 f(p,q)\leq\frac35.
\]
The supremum \(3/5\) occurs at the weak endpoint \(C_0\) and is
approached by strong interior points.
\end{proposition}

\begin{proof}
Every vertex in \eqref{eq:mixed-hull} satisfies \(x+y\geq1/2\);
equality occurs only at \(A_0\), which fails the other Region IV
condition.  Hence the hull does not enter Regions IV--V.

Suppose first that \(x\leq1/3\), so \(p\geq3\) and Regions I--II are
the relevant ones.  The affine inequality
\begin{equation}
 \label{eq:affine-six}
 12x+11y\geq6
\end{equation}
holds at all three vertices and hence throughout the hull.  In the
positive-denominator region of \eqref{eq:f12},
\eqref{eq:affine-six} is equivalent to
\[
 \frac{2y}{1-2x-y}\geq\frac{12}{5}.
\]
Thus \(f\leq3/5\).

Now suppose \(1/3<x<1/2\).  The hull lies above its edge
\(A_0C_0\):
\begin{equation}
 \label{eq:lower-edge}
 y\geq\frac6{11}(1-2x).
\end{equation}
Only Region III is relevant.  For fixed \(x\), the quotient subtracted
in \eqref{eq:f3} increases with \(y\), so \(f\) is maximised at equality
in \eqref{eq:lower-edge}.  Put \(z=1-2x\in(0,1/3)\).  Substitution gives
\[
 f
 =
 3-\frac{24(2-3z)}{11-3z}.
\]
This is strictly increasing in \(z\), and its limit at \(z=1/3\),
the vertex \(C_0\), equals \(3/5\).
\end{proof}

\section{The terminal dimension pincer}

Let \(\EE\) be a weak-star limit of
\(|u(s_j)|^2\,dx\).  The uniform spatial tightness proved in the
preceding temporal theorem gives
\begin{equation}
 \label{eq:mass-tight}
 \EE(\R^3)=E_-.
\end{equation}
Weak continuity gives
\begin{equation}
 \label{eq:weak-terminal}
 u(s_j)\weakto u_*
 \quad\hbox{in }L^2.
\end{equation}
Convex lower semicontinuity on nonnegative compactly supported tests
therefore gives the positive measure
\begin{equation}
 \label{eq:defect-measure}
 \vartheta:=\EE-|u_*|^2\,dx\geq0.
\end{equation}
Equations \eqref{eq:energy-defect} and \eqref{eq:mass-tight} give
\begin{equation}
 \label{eq:defect-mass}
 \vartheta(\R^3)
 =
 E_--\|u_*\|_2^2
 \geq d_0>0.
\end{equation}

Define the terminal singular set
\[
 \sigma
 :=
 \{x:(x,T^*)\hbox{ is a singular point of the continuation}\}.
\]

\begin{lemma}[Support of the defect]
\label{lem:support}
\[
 \supp\vartheta\subset\sigma.
\]
\end{lemma}

\begin{proof}
Let \(K\Subset\R^3\setminus\sigma\).  Finitely many regular
parabolic neighbourhoods cover \(K\times\{T^*\}\).  Local regularity
in smaller cylinders gives strong \(L^2(K)\) convergence of \(u(t)\)
to the weak trace \(u_*\) as \(t\uparrow T^*\).  Therefore
\[
 \EE\lfloor_K=|u_*|^2\,dx\lfloor_K.
\]
Exhaust the regular open set by compact subsets.
\end{proof}

\begin{theorem}[Terminal dimension pincer]
\label{thm:pincer}
Every Leray--Hopf continuation in the conditional energy-efficient
exact \(q=4\) branch satisfies
\begin{equation}
 \label{eq:pincer}
 \boxed{
 \frac35\leq\dimH\sigma\leq1,
 \qquad
 \HH^1(\sigma)=0.
 }
\end{equation}
\end{theorem}

\begin{proof}
The temporal five-power result and
Proposition~\ref{prop:three-fifths} give the
Leslie--Shvydkoy concentration-dimension bound
\[
 D(\EE)\geq\frac35.
\]
Consequently, \(\EE(S)=0\) for every Borel set with
\(\dimH S<3/5\).  Since \(0\leq\vartheta\leq\EE\), the same holds for
\(\vartheta\).

If \(\dimH\sigma<3/5\), Lemma~\ref{lem:support} and
\eqref{eq:defect-mass} would give
\[
 0<\vartheta(\R^3)=\vartheta(\sigma)=0,
\]
a contradiction.  This proves the lower bound.  The
Caffarelli--Kohn--Nirenberg theorem gives \(\HH^1(\sigma)=0\), hence
the upper bound.
\end{proof}

\begin{remark}
The strict interval from \(3/5\) to \(1\) is not closed by the current
information.  Theorem~\ref{thm:pincer} is a necessary property of the
conditional survivor, not a proof that such a survivor exists.
\end{remark}

\section{Pressure audit}

For the whole-space pressure representative
\[
 \pi=\mathcal R_a\mathcal R_b(u_au_b),
\]
Lorentz Calder\'on--Zygmund bounds and
Lemma~\ref{lem:endpoint} imply
\begin{equation}
 \label{eq:pressure-endpoint}
 \pi\in
 L_t^{11/4,\infty}L_x^{3/2,\infty}.
\end{equation}
Its pressure scaling index is
\begin{equation}
 \label{eq:pressure-gamma}
 \frac2{11/4}+\frac3{3/2}
 =
 2+\frac8{11}.
\end{equation}
The ordinary energy class generates the pressure line
\[
 \frac2r+\frac3p=3.
\]
Every direct interpolation between that line and
\eqref{eq:pressure-endpoint} therefore satisfies
\begin{equation}
 \label{eq:gamma-lower}
 \frac2r+\frac3p=2+\gamma,
 \qquad
 \gamma\geq\frac8{11}.
\end{equation}

Barker and Wang's supercritical theorem assumes
\[
 0<\gamma<\frac12,
 \qquad
 \delta>\frac{3\gamma}{2+2\gamma},
\]
as well as its further \(p,r,\delta,\gamma\) restrictions, and then
concludes \(\HH^{2\delta}(\sigma)=0\).  The direct q4 hull
\eqref{eq:gamma-lower} is already outside the theorem.

Moreover, entering the pincer-closing range \(2\delta<3/5\) would
require, before the further restrictions,
\[
 \frac{3\gamma}{1+\gamma}<\frac35
 \quad\Longleftrightarrow\quad
 \gamma<\frac14.
\]
Thus the direct pressure interpolation misses this necessary closing
range by at least
\begin{equation}
 \label{eq:pressure-gap}
 \boxed{
 \frac8{11}-\frac14=\frac{21}{44}.
 }
\end{equation}
This does not preclude a structural pressure-cancellation or
pressure-oscillation gain.

\section{Clock-band correlation or super-clock escape}

Set
\[
 w_j:=u(s_j)-u_*.
\]
Then \(w_j\weakto0\) in \(L^2\).  Since the high-frequency tail of the
fixed \(u_*\) vanishes, \eqref{eq:old-floor} gives, after harmless
fixed cutoff changes,
\begin{equation}
 \label{eq:w-tail}
 \|\Pi_{>\Lambda_j}w_j\|_2\geq d_1>0.
\end{equation}
Here \(\Pi\) is the sharp orthogonal Fourier projection.  It dominates
the smooth high-pass multiplier used to obtain
\eqref{eq:old-floor}.

\begin{theorem}[Exhaustive spectral alternative]
\label{thm:spectral}
At least one of the following holds; mutual exclusivity is not asserted.

\medskip
\noindent
\textbf{A. Bounded clock band.}
There are fixed \(A>1\), \(\eta>0\), and infinitely many \(j\) such
that
\begin{equation}
 \label{eq:clock-band}
 \left\|
 \Pi_{\Lambda_j<|\xi|\leq A\Lambda_j}w_j
 \right\|_2^2
 \geq\eta.
\end{equation}
For a further subsequence \(j_k\), define
\[
 g_k
 :=
 \Pi_{\Lambda_{j_k}<|\xi|\leq A\Lambda_{j_k}}w_{j_k},
 \qquad
 I_k=[s_{j_k},s_{j_{k+1}}].
\]
The intervals are disjoint and, for all large \(k\),
\begin{equation}
 \label{eq:signed-work}
 \boxed{
 \int_{I_k}\!\!\int_{\R^3}
 u\otimes u:\nabla g_k\,dx\,dt
 \leq-\frac{\eta}{4}.
 }
\end{equation}
The corresponding viscous work tends to zero:
\begin{equation}
 \label{eq:viscous-zero}
 \left|
 \nu\int_{I_k}\!\!\int_{\R^3}
 \nabla u:\nabla g_k\,dx\,dt
 \right|\longrightarrow0.
\end{equation}

\medskip
\noindent
\textbf{B. Super-clock escape.}
If Alternative A fails, there are \(A_k\to\infty\) and a subsequence
\(j_k\) such that
\begin{equation}
 \label{eq:superclock-tail}
 \left\|
 \Pi_{>A_k\Lambda_{j_k}}u(s_{j_k})
 \right\|_2
 \geq\eta_0>0.
\end{equation}
Consequently,
\begin{equation}
 \label{eq:superclock-enstrophy}
 \boxed{
 \|\nabla u(s_{j_k})\|_2^2
 \gtrsim
 A_k^2\Lambda_{j_k}^2.
 }
\end{equation}
\end{theorem}

\begin{proof}
Assume first that \eqref{eq:clock-band} holds.  For each fixed
\(g_k\), weak convergence \(w_j\weakto0\) permits the next index to be
chosen so far out that
\[
 |\langle w_{j_{k+1}},g_k\rangle|
 \leq\frac{\eta}{4}.
\]
Because \(g_k\) is an orthogonal projection of \(w_{j_k}\),
\begin{align}
 \label{eq:correlation-drop}
 \left\langle
 u(s_{j_{k+1}})-u(s_{j_k}),g_k
 \right\rangle
 &=
 \langle w_{j_{k+1}},g_k\rangle-\|g_k\|_2^2\\
 &\leq-\frac{3\eta}{4}.\nonumber
\end{align}
Pair the smooth equation on \(I_k\) with the fixed divergence-free
field \(g_k\):
\begin{align}
 \label{eq:work-identity}
 \left\langle
 u(s_{j_{k+1}})-u(s_{j_k}),g_k
 \right\rangle
 ={}&
 \int_{I_k}\!\!\int u\otimes u:\nabla g_k\\
 &-
 \nu\int_{I_k}\!\!\int\nabla u:\nabla g_k.\nonumber
\end{align}

The exact record tail gives
\begin{equation}
 \label{eq:two-tail-bounds}
 T^*-s_j\lesssim\frac{m_j^{-11/2}}j,
 \qquad
 \int_{s_j}^{T^*}M(t)^2\,dt
 \lesssim\frac{m_j^{-7/2}}j.
\end{equation}
Band-limited Lorentz Bernstein gives
\[
 \|\nabla g_k\|_{L^{3,1}}
 \lesssim
 (A\Lambda_{j_k})^{3/2}\|g_k\|_2.
\]
Lorentz H\"older and \eqref{eq:two-tail-bounds} therefore yield
\begin{align}
 \label{eq:critical-capacity}
 \int_{s_j}^{T^*}
 \left|\int u\otimes u:\nabla g_j\right|dt
 &\lesssim
 (A\Lambda_j)^{3/2}
 \int_{s_j}^{T^*}M(t)^2\,dt\\
 &\lesssim A^{3/2},\nonumber
\end{align}
because
\(\Lambda_j^{3/2}\asymp m_j^{7/2}j\).
The nonlinear capacity is exactly order one per event.

Set
\[
 D_j^\nu
 :=
 \nu\int_{s_j}^{T^*}\|\nabla u(t)\|_2^2\,dt.
\]
Absolute continuity gives \(D_j^\nu\to0\).  Cauchy--Schwarz,
\eqref{eq:two-tail-bounds}, and
\(\|\nabla g_j\|_2\leq A\Lambda_j\|g_j\|_2\) give
\begin{align}
 \label{eq:viscous-tail}
 \left|
 \nu\int_{s_j}^{T^*}\!\!\int\nabla u:\nabla g_j
 \right|
 &\lesssim
 A\Lambda_j(T^*-s_j)^{1/2}(D_j^\nu)^{1/2}\\
 &\lesssim
 A\,m_j^{-5/12}j^{1/6}(D_j^\nu)^{1/2}
 \longrightarrow0.\nonumber
\end{align}
Equations \eqref{eq:correlation-drop}--\eqref{eq:viscous-tail}
give \eqref{eq:signed-work}--\eqref{eq:viscous-zero}.

If Alternative A fails, then for each fixed \(A>1\),
\[
 \left\|
 \Pi_{\Lambda_j<|\xi|\leq A\Lambda_j}w_j
 \right\|_2\longrightarrow0.
\]
Combine this with \eqref{eq:w-tail} and choose
\(A_k\to\infty\) diagonally to obtain
\[
 \|\Pi_{>A_k\Lambda_{j_k}}w_{j_k}\|_2
 \geq d_1/2.
\]
Since \(A_k\Lambda_{j_k}\to\infty\),
\[
 \|\Pi_{>A_k\Lambda_{j_k}}u_*\|_2\longrightarrow0.
\]
The triangle inequality proves \eqref{eq:superclock-tail}, and
Plancherel proves \eqref{eq:superclock-enstrophy}.
\end{proof}

\begin{remark}[Why Alternative A does not yet sum]
The work in \eqref{eq:signed-work} has a fixed sign on disjoint
intervals, but every test \(g_k\) changes with the event.  The capacity
\eqref{eq:critical-capacity} is order one on every interval, so the
current norms do not furnish one finite common budget for the
unweighted sum.  A transported adjoint or a genuine square-function
estimate would have to exploit more than separate Lorentz duality.
\end{remark}

\begin{remark}[Why Alternative B does not yet dissipate]
Equation \eqref{eq:superclock-enstrophy} is pointwise.  No lower
residence time is known at the frequency
\(A_k\Lambda_{j_k}\), and the factor \(A_k\) may be accompanied by an
arbitrarily shorter interval.  Thus the enstrophy amplification is not
an integrated dissipation contradiction.
\end{remark}

\section{What changed in this round}

\subsection*{Formal findings, subject to external review}

\begin{enumerate}
\item The exact q4 record clock reaches weak
\(L_t^{11/2}\) and every Lorentz refinement
\(L_t^{11/2,r}\), \(r>11/2\).
\item Even under the additional favourable surrogate assumption
\(L_t^{11/2,\infty}B^0_{3,\infty}\), the Cheskidov--Luo Besov threshold
is missed throughout the complete direct interpolation hull.
\item The Leslie--Shvydkoy threshold has exact hull supremum \(3/5\).
\item The nonzero terminal energy defect is supported on the terminal
singular set and forces the pincer \eqref{eq:pincer}.
\item Direct pressure interpolation has
\(\gamma\geq8/11\) and misses the Barker--Wang
dimension-closing necessary range \(\gamma<1/4\) by \(21/44\).
\item The terminal defect obeys the exhaustive clock-band or
super-clock alternative of Theorem~\ref{thm:spectral}.
\end{enumerate}

\subsection*{Closed shortcuts}

\begin{enumerate}
\item Endpoint weak time integrability alone does not close energy
equality: the missing spatial Besov smoothness is explicit in
\eqref{eq:four-gaps}.
\item Reoptimising all direct energy-class mixed norms cannot raise the
Leslie--Shvydkoy threshold above \(3/5\).
\item Squaring the q4 velocity bound into a pressure bound does not
enter the Barker--Wang parameter region.
\item Orthogonal or signed defect correlations alone are not additive,
because their tests change and the nonlinear capacity is exactly
critical.
\item A larger pointwise enstrophy floor alone is not an integrated
charge without residence.
\end{enumerate}

\subsection*{Things left to prove}

\begin{enumerate}
\item Sum or cancel the event-dependent nonlinear works in
\eqref{eq:signed-work} against one finite same-trajectory budget.
\item Turn super-clock escape into a positive residence or dissipation
lower bound.
\item Derive a structural pressure cancellation or oscillation gain
that crosses the direct \(21/44\) deficit.
\item Prove strong left \(L^2\) continuity or energy equality at
\(T^*\).
\item Control the divergent-normalised-energy branch and schedules at
or below the temporal five-power barrier.
\item Prove one complete Clay alternative for arbitrary admissible
data.
\end{enumerate}

\begin{conjecture}[The pincer is empty]
No first singular time arising from smooth finite-energy data supports
a nonzero anomalous energy defect on a terminal singular set of
Hausdorff dimension in \([3/5,1]\).
\end{conjecture}

Theorem~\ref{thm:pincer} proves that every energy-efficient exact q4
survivor would have precisely such a defect.  It does not prove the
conjecture, a global regularity theorem, a breakdown construction, or
any Clay alternative.

\begin{thebibliography}{9}
\raggedright

\bibitem{BW23}
T.~Barker and W.~Wang,
\emph{Estimates of the singular set for the Navier--Stokes equations
with supercritical assumptions on the pressure},
J. Differential Equations \textbf{365} (2023), 379--407.
\href{https://doi.org/10.1016/j.jde.2023.04.007}
{doi:10.1016/j.jde.2023.04.007};
arXiv:2111.15444v2.

\bibitem{CKN82}
L.~Caffarelli, R.~Kohn, and L.~Nirenberg,
\emph{Partial regularity of suitable weak solutions of the
Navier--Stokes equations},
Comm. Pure Appl. Math. \textbf{35} (1982), 771--831.
\href{https://doi.org/10.1002/cpa.3160350604}
{doi:10.1002/cpa.3160350604}.

\bibitem{CL20}
A.~Cheskidov and X.~Luo,
\emph{Energy equality for the Navier--Stokes equations in weak-in-time
Onsager spaces},
Nonlinearity \textbf{33} (2020), 1388--1403.
\href{https://doi.org/10.1088/1361-6544/ab60d3}
{doi:10.1088/1361-6544/ab60d3};
arXiv:1802.05785v2.

\bibitem{LS18}
T.~M. Leslie and R.~Shvydkoy,
\emph{The energy measure for the Euler and Navier--Stokes equations},
Arch. Ration. Mech. Anal. \textbf{230} (2018), 459--492.
\href{https://doi.org/10.1007/s00205-018-1250-4}
{doi:10.1007/s00205-018-1250-4};
arXiv:1705.04420v4.

\end{thebibliography}

\end{document}
